\sect{Large Language Models and Function Calling}{functioncalls}

LLMs are large-scale transformer-based architectures trained on extensive text corpora to predict and generate language in context. Beyond text generation, they exhibit emergent reasoning, abstraction, and semantic parsing abilities, allowing them to interpret unstructured NL input and express it in structured, machine-interpretable formats \cite{10.1145/3744746}. These capabilities make LLMs suitable not only for linguistic tasks but also for integration with external computational systems where deterministic execution and contextual reasoning must coexist.

Function calling extends these abilities by providing a structured interface through which an LLM can invoke predefined software functions or tools based on user intent. In this paradigm, the model produces a JSON or schema-constrained output specifying which function to call and what parameters to pass \cite{wang2025function}. The process effectively bridges free-form NL and formal command execution by transforming intent recognition into parameterized function invocation. This structured interaction enables LLMs to operate within rule-bound environments such as control systems, where precision and safety are critical \cite{10.1145/3744746}.

In practice, function calling can be viewed as a three-stage reasoning cycle: intent recognition, tool selection, and argument formulation. The model first interprets the user’s goal and retrieves the appropriate function from a registered set of tools; it then extracts and formats the required parameters to meet the target schema. Validation layers or mapping strategies ensure syntactic and semantic compatibility between the model output and the execution environment, reducing ambiguity and minimizing failure rates \cite{wang2025function}. This structure allows the LLM to act as a reasoning and coordination layer while delegating deterministic control to external systems.

When extended to autonomous agents, function calling becomes the core mechanism that enables multi-step planning and interaction with real or simulated environments. An AI agent equipped with such capabilities can reason about sequences of actions, call external tools to perform them, and evaluate the resulting state before proceeding. The agent’s behavior is typically constrained by guardrails, explicit or learned safety filters, that restrict tool use, monitor policy compliance, and validate outputs to prevent unintended or unsafe operations \cite{openai2023practicalagents}. Through this architecture, agents achieve both autonomy and controllability, a combination that is essential for integration in physical domains.

Recent frameworks implement this paradigm with varying degrees of abstraction and transparency. LangChain provides a modular orchestration layer that manages prompt templates, memory, and tool invocation pipelines, making it suitable for building composite reasoning workflows but often at the cost of interpretability and efficiency. Llama, in contrast, specializes in structured data retrieval and context synthesis, complementing agents that rely on precise grounding in domain-specific knowledge. SmolAgent, developed by Hugging Face, adopts a code-generation-centric approach in which the model directly produces and executes Python instructions, offering fine-grained control over tool usage and state management. This explicit code execution model, though less abstracted, yields higher transparency and reproducibility, features particularly advantageous in scientific and robotic environments where deterministic behavior and verifiable execution traces are required \cite{towardsai2024comparison}.

Recent advances in lightweight agent frameworks have shifted focus from large orchestration stacks toward transparent, code-centric execution models. The \textit{smolagents} library by Hugging~Face exemplifies this transition by enabling agents that reason through explicit code generation rather than predefined toolchains \cite{huggingface_smolagents_good_agents}. In this architecture, the model produces self-contained Python functions that directly interact with local or remote tools, execute environment queries, and process outputs in an interpretable manner. This approach contrasts with high-level frameworks that conceal execution logic behind abstraction layers, offering instead full observability of reasoning and action steps. Empirical demonstrations further indicate that such minimalistic agents can achieve competitive performance in complex multi-step reasoning tasks while maintaining high reproducibility and reduced dependency overheads \cite{datacamp_smolagents_guide}. These characteristics position code-executing agents as a promising foundation for safety-critical or physically grounded systems, where deterministic traceability and modular integration outweigh large-scale orchestration complexity.


These developments have led to practical systems where LLM-based agents can translate high-level NL objectives into executable functions operating on robots, sensors, and analytical instruments \cite{huang2022language}. In such settings, abstract instructions such as “measure the OCP of the first sample” or “bring the third specimen to the user area” can be transformed into structured calls like \texttt{measure\_ocp(i=1)} or \texttt{collect\_sample\_from\_user(i=3)}. This direct linkage between NL intent and physical execution forms the conceptual foundation of agentic control architectures for collaborative autonomous environments.
